%%% Кодировки и шрифты %%%
\ifxetexorluatex
    % Язык по-умолчанию русский с поддержкой приятных команд пакета babel
    \setmainlanguage[babelshorthands=true]{russian}
    % Языки библиографии: langid выбирает локальные сокращения и переносы.
    \setotherlanguages{english,catalan,french,german,italian,latin,portuguese,spanish}

    % Проверка существования шрифтов. Недоступна в pdflatex
    \ifnumequal{\value{fontfamily}}{1}{
        \IfFontExistsTF{Times New Roman}{}{\setcounter{fontfamily}{0}}
    }{}
    \ifnumequal{\value{fontfamily}}{2}{
        \IfFontExistsTF{LiberationSerif}{}{\setcounter{fontfamily}{0}}
    }{}

    \ifnumequal{\value{fontfamily}}{0}{% Семейство шрифтов CMU. Используется как fallback
        \setmonofont[
          BoldFont={cmuntb.otf},
          ItalicFont={cmunit.otf},
          BoldItalicFont={cmuntx.otf}
        ]{cmuntt.otf} % {CMU Typewriter Text} % моноширинный шрифт
        \newfontfamily\cyrillicfonttt[
          BoldFont={cmuntb.otf},
          ItalicFont={cmunit.otf},
          BoldItalicFont={cmuntx.otf}
        ]{cmuntt.otf} % {CMU Typewriter Text} % моноширинный шрифт для кириллицы
        \defaultfontfeatures{Ligatures=TeX}   % стандартные лигатуры TeX, замены нескольких дефисов на тире и т. п. Настройки моноширинного шрифта должны идти до этой строки, чтобы при врезках кода программ в коде не применялись лигатуры и замены дефисов
        \setmainfont[
          SlantedFont={cmunsl.otf},
          BoldSlantedFont={cmunbl.otf},
          BoldFont={cmunbx.otf},
          ItalicFont={cmunti.otf},
          BoldItalicFont={cmunbi.otf}
        ]{cmunrm.otf} % {CMU Serif}           % Шрифт с засечками
        \newfontfamily\cyrillicfont[
          SlantedFont={cmunsl.otf},
          BoldSlantedFont={cmunbl.otf},
          BoldFont={cmunbx.otf},
          ItalicFont={cmunti.otf},
          BoldItalicFont={cmunbi.otf}
        ]{cmunrm.otf}   % {CMU Serif}         % Шрифт с засечками для кириллицы
        \setsansfont[
          BoldFont={cmunsx.otf},
          ItalicFont={cmunsi.otf},
          BoldItalicFont={cmunso.otf}
        ]{cmunss.otf} % {CMU Sans Serif}      % Шрифт без засечек
        \newfontfamily\cyrillicfontsf[
          BoldFont={cmunsx.otf},
          ItalicFont={cmunsi.otf},
          BoldItalicFont={cmunso.otf}
        ]{cmunss.otf} % {CMU Sans Serif}      % Шрифт без засечек для кириллицы
    }

    \ifnumequal{\value{fontfamily}}{1}{                    % Семейство MS шрифтов
        \setmonofont{Times New Roman}                      % единый шрифт, включая листинги
        \newfontfamily\cyrillicfonttt{Times New Roman}     % единый шрифт для кириллицы
        \defaultfontfeatures{Ligatures=TeX}                % стандартные лигатуры TeX, замены нескольких дефисов на тире и т. п. Настройки моноширинного шрифта должны идти до этой строки, чтобы при врезках кода программ в коде не применялись лигатуры и замены дефисов
        \setmainfont{Times New Roman}                      % Шрифт с засечками
        \newfontfamily\cyrillicfont{Times New Roman}       % Шрифт с засечками для кириллицы
        \setsansfont{Times New Roman}                      % единый шрифт во всём документе
        \newfontfamily\cyrillicfontsf{Times New Roman}     % единый шрифт для кириллицы
        \ifxetex
            \setmathsfont(Digits,Latin,Greek){Times New Roman}
            \setmathrm{Times New Roman}
            \setmathsf{Times New Roman}
            \setmathtt{Times New Roman}
        \fi
    }

    \ifnumequal{\value{fontfamily}}{2}{                    % Семейство шрифтов Liberation (https://pagure.io/liberation-fonts)
        \setmonofont{LiberationMono}[Scale=0.87] % моноширинный шрифт
        \newfontfamily\cyrillicfonttt{LiberationMono}[     % моноширинный шрифт для кириллицы
            Scale=0.87]
        \defaultfontfeatures{Ligatures=TeX}                % стандартные лигатуры TeX, замены нескольких дефисов на тире и т. п. Настройки моноширинного шрифта должны идти до этой строки, чтобы при врезках кода программ в коде не применялись лигатуры и замены дефисов
        \setmainfont{LiberationSerif}                      % Шрифт с засечками
        \newfontfamily\cyrillicfont{LiberationSerif}       % Шрифт с засечками для кириллицы
        \setsansfont{LiberationSans}                       % Шрифт без засечек
        \newfontfamily\cyrillicfontsf{LiberationSans}      % Шрифт без засечек для кириллицы
    }

\else
    \ifnumequal{\value{usealtfont}}{1}{% Используется pscyr, при наличии
        \IfFileExists{pscyr.sty}{\renewcommand{\rmdefault}{ftm}}{}
    }{}
\fi
